\begin{theorem}[Задача о шести знакомых]
Среди любых шести человек найдутся трое попарно знакомых или трое попарно
незнакомых. На языке графов: любой граф \(G\) на 6 вершинах содержит \(K_3\)
или независимое множество размера 3 (т.\,е. \(K_3\subseteq\overline G\)).
\end{theorem}
\begin{proof}
Возьмём вершину \(x\): \(d_G(x)+d_{\overline G}(x)=5\), поэтому одна из степеней
не меньше 3. Переходя при необходимости к \(\overline G\) (утверждение симметрично),
считаем \(d_G(x)\ge3\) и \(a,b,c\in N(x)\). Если какие-то две из \(a,b,c\) смежны,
вместе с \(x\) они образуют треугольник. Иначе \(\{a,b,c\}\) независимо.
\end{proof}
\begin{graphillustration}{Доказательство для шести вершин: у $x$ три соседа $a,b,c$. Любое ребро между ними (пунктир) замыкает треугольник с $x$; если рёбер нет — $\{a,b,c\}$ независимо.}{fig:07-six}
\begin{tikzpicture}[font=\small]
\node[gi selected] (x) at (0,1.5) {$x$};
\node[gi vertex] (a) at (-1.5,0) {$a$};
\node[gi vertex] (b) at (0,0) {$b$};
\node[gi vertex] (c) at (1.5,0) {$c$};
\node[gi vertex] (d) at (2.8,1.2) {};
\node[gi vertex] (e) at (3.3,0.2) {};
\draw[gi c1] (x) -- (a); \draw[gi c1] (x) -- (b); \draw[gi c1] (x) -- (c);
\draw[gi c2] (x) -- (d); \draw[gi c2] (x) -- (e);
\draw[gi edge,densely dotted] (a) -- (b);
\draw[gi edge,densely dotted] (b) -- (c);
\draw[gi edge,densely dotted] (a) to[bend right=30] (c);
\node[gi note,anchor=west] at (4.2,0.8) {$d_G(x)\ge3$ (сплошные);\\штриховые — рёбра $\overline G$};
\end{tikzpicture}
\end{graphillustration}


\begin{definition}[Число Рамсея]
\(R(m,n)\) — наименьшее \(k\in\N\) такое, что любой граф на \(k\) вершинах содержит
клику \(K_m\) или независимое множество размера \(n\) (\(\overline{K_n}\)).
Эквивалентно: при любой раскраске рёбер \(K_k\) в два цвета есть клика \(K_m\)
цвета 1 или клика \(K_n\) цвета 2.
\end{definition}
\begin{lemma}[Смысл определения]
Пусть \(R(m,n)=k\). Тогда (1) любой граф на \(\ge k\) вершинах содержит \(K_m\) или
независимое множество размера \(n\); (2) для любого \(t<k\) существует граф на \(t\)
вершинах без \(K_m\) и без независимого множества размера \(n\).
\end{lemma}
\begin{proof}
(1) Если \(|V|\ge k\), выберем \(V'\subseteq V\), \(|V'|=k\), и применим определение к
\(G[V']\); найденная клика или независимое множество годятся и в \(G\).
(2) При \(t=k-1\) пример существует по минимальности \(k\); при меньших \(t\) возьмём
его индуцированный подграф на \(t\) вершинах.
\end{proof}

\begin{tcolorbox}[colback=cbGrayBack,colframe=cbGray!55]
\textbf{Простейшие значения и свойства.}
\begin{itemize}\tightlist
\item \(R(m,n)=R(n,m)\): переход к дополнению меняет клики и независимые множества местами.
\item \(R(1,n)=1\): одна вершина — уже \(K_1\).
\item \(R(2,n)=n\): граф на \(n\) вершинах либо имеет ребро (\(K_2\)), либо пуст
  (независим целиком); пустой граф на \(n-1\) вершинах — контрпример.
\item Монотонность: \(R(m,n)\le R(m',n')\) при \(m\le m'\), \(n\le n'\).
\end{itemize}
\end{tcolorbox}

\begin{cor}
\(R(3,3)=6\).
\end{cor}
\begin{proof}
\(R(3,3)\le6\) по теореме о шести знакомых. Цикл \(C_5\) не содержит треугольника,
а его дополнение — снова \(C_5\) (пентаграмма), т.\,е. нет и независимого множества
размера 3. Поэтому \(R(3,3)>5\).
\end{proof}
\begin{graphillustration}{Раскраска $K_5$ без одноцветного треугольника: цвет 1 — цикл $C_5$, цвет 2 — пентаграмма (тоже $C_5$). Поэтому $R(3,3)>5$.}{fig:07-c5}
\begin{tikzpicture}[font=\small]
\foreach \i in {0,...,4} \node[gi vertex] (v\i) at (90+72*\i:1.5) {$\i$};
\foreach \i [evaluate=\i as \j using {int(mod(\i+1,5))}] in {0,...,4} \draw[gi c1] (v\i) -- (v\j);
\foreach \i [evaluate=\i as \j using {int(mod(\i+2,5))}] in {0,...,4} \draw[gi c2] (v\i) -- (v\j);
\node[gi note,anchor=west] at (2.3,0.4) {сплошные: $G=C_5$\\штриховые: $\overline G\cong C_5$};
\end{tikzpicture}
\end{graphillustration}

Конечность \(R(m,n)\) для всех \(m,n\) будет доказана в билете~\ref{ticket:8}.
Известно немного точных значений: \(R(3,4)=9\), \(R(3,5)=14\), \(R(4,4)=18\);
уже \(R(5,5)\) неизвестно (\(43\le R(5,5)\le46\)).
